Campus sustainability platforms usually centralize telemetry and display historical energy, water, waste, and environmental indicators. They provide visibility but rarely combine local multimodal inference, uncertainty-aware prediction, verified carbon accounting, and safe real-time control. Their cloud dependence also creates latency, connectivity, bandwidth, privacy, and operational-resilience limitations. This proposal specifies \system{}, an edge-AI platform that executes building forecasting, fault detection, and carbon-aware supervisory optimization on a Qualcomm Innovators Development Kit (QIDK) Android device. The device consumes electricity, photovoltaic, HVAC, indoor-air-quality, occupancy, water, waste, weather, tariff, and grid-carbon data through a segmented protocol gateway. It cannot replace certified building controllers or life-safety systems. The principal control method is constrained model predictive control (MPC), with deterministic limits and fallback sequences independent of learned models. Recent time-series foundation models and graph models are evaluated as challengers, not presumed superior to seasonal and gradient-boosted baselines. Models are quantized and compiled with Qualcomm AI Runtime and QNN for CPU, GPU, and Hexagon NPU comparison. A staged replay, digital-twin, hardware-in-the-loop, shadow, advisory, and randomized field evaluation measures adjusted energy, operational carbon, peak demand, comfort, indoor air quality, availability, privacy, inference latency, and edge-compute energy. Savings follow an IPMVP-style counterfactual and report location- and market-based emissions separately. The central hypothesis is that local, uncertainty-gated supervisory intelligence can reduce operational energy and emissions while preserving occupant service, auditability, privacy, and safe building operation.
